\documentclass[11pt]{article}

\usepackage{multicol}
\usepackage{times}
\usepackage{graphicx}
\usepackage{epsfig}

\textheight 9in
\textwidth 6.5in
\oddsidemargin 0in
\evensidemargin 0in
\topmargin -.5in
\pagestyle{empty}
\newcommand{\setspace}[1]{\renewcommand{\baselinestretch}{#1}\large\normalsize}
%\setspace{1.1}
\renewcommand{\textfraction}{0.1}
\renewcommand{\topfraction}{0.9}
\newcommand{\tods}{{\em TODS\/}}

\begin{document}
\begin{center}
{\Large\bf Developments at ACM {\em TODS}}\\[2ex]
{\large\it  Richard Snodgrass}\\[1ex]
{\tt rts@cs.arizona.edu}\\[3ex]
\end{center}

The March 2005 issue of \tods\ has eight papers invited from the SIGMOD and
PODS'2003 conferences. These papers are significantly extended versions of
the conference papers, allowing the authors to refine and elaborate without
the strictures of a twelve-page limit.

The first four papers in that issue were invited from the SIGMOD conference;
the last four papers were invited from the PODS conference. Each went
through the normal rigorous review process.

Alon Halevy was Program Chair for SIGMOD'2003; Tova Milo was Program Chair
for PODS'2003. It is interesting that the first author of the first paper in
this special issue is Tova. I should emphasize that while Tova helped
select the PODS papers to invite, she had no involvement in selecting the
SIGMOD papers to invite. It just turned out that one of the four SIGMOD
papers selected was co-authored by her.

The June 2005 issue is almost
complete; it will have at least seven papers (available now on the Upcoming
Issues page\footnote{{\tt http://www.acm.org/tods/Upcoming.html}}). \tods\
continues to grow: the first two issues of 2005 contain more papers than
the first two issues of 2004, and more than the number of papers that appeared
in all of 2002. The full story is on the \tods\ web site\footnote{\tt
  http://www.acm.org/tods/TurnaroundTime.html}. 

\vspace{2em}
I have appointed six new Associate Editors, bringing the complete Editorial
Board to nineteen.

\begin{description}
\item[Jan Chomicki's] research interests are in logical foundations of
databases.  Specific topics include: database integrity, data integration,
data models, and query languages. His current projects involve query
answering in inconsistent databases and preference queries.

\item[Heikki Mannila] works in data mining, algorithms and bioinformatics.

\item[Raghu Ramakrishnan's] current research is in two broad areas. In the EDAM
project, he is working on data mining problems with driving applications in
environmental monitoring, physics simulations, and e-commerce.  In the
CICADA project, he is working together with researchers from Microsoft
Research on extending SQL to allow applications to specify when (potentially
out of date) copies of data can be used.

\item[Arnie Rosenthal's] research is in the areas of data security and data
sharing. In particular, he tries to align the technologies with feasible
practices for data-owning organizations.

\item[Sunita Sarawagi's] research interests span several fields including 
databases, data mining, machine learning and statistics.  Currently she 
is investigating the deployment of learning-based techniques for solving 
various data integration and cleaning tasks.

\item[Dan Suciu] works on applying formal theory to novel and difficult data 
management tasks.  His past work was on various aspects of managing 
semistructured data, including query languages, compression, query 
processing and type inference, while his recent work focused on data 
security and on querying unreliable and inconsistent data sources.
\end{description}

\noindent
All six are internationally-known scholars in the field of
database systems and are well known also to the database community
through their past service. In addition, they hail from three different
continents.

I'm gratified that they are willing to help \tods\ continue to improve.

\end{document}
